\documentclass{article}

\PassOptionsToPackage{numbers}{natbib}
\usepackage[preprint]{neurips_2026}

\usepackage[utf8]{inputenc}
\usepackage[T1]{fontenc}
\usepackage{hyperref}
\usepackage{url}
\usepackage{booktabs}
\usepackage{amsfonts}
\usepackage{amsmath}
\usepackage{amssymb}
\usepackage{nicefrac}
\usepackage{microtype}
\usepackage{xcolor}
\usepackage{graphicx}
\usepackage{enumitem}

\title{Per-Unit System: A Reference for Power-Flow Computations}

\author{Changhun Kim}

\begin{document}
\maketitle

%==============================================================================
\section{The Four Base Quantities}
%==============================================================================

In the per-unit (p.u.) system, all electrical quantities are expressed as
dimensionless ratios relative to a chosen \emph{base}.  There are four base
quantities, but only \textbf{two are chosen freely}:

\begin{itemize}[noitemsep]
  \item $S_{\rm base}$ [VA] — one value for the entire system (commonly
    $100\,\mathrm{MVA}$ for transmission, $1\,\mathrm{MVA}$ for distribution).
  \item $V_{\rm base}$ [V] — one value \emph{per bus}, which can differ
    across buses in a multi-voltage network.
\end{itemize}

The remaining two bases are \textbf{derived automatically}:

\begin{equation}
  I_{\rm base} = \frac{S_{\rm base}}{V_{\rm base}},
  \qquad
  Z_{\rm base} = \frac{V_{\rm base}^2}{S_{\rm base}},
  \qquad
  Y_{\rm base} = \frac{1}{Z_{\rm base}} = \frac{S_{\rm base}}{V_{\rm base}^2}
  \label{eq:bases}
\end{equation}

The key relationship to remember: $Y_{\rm base}$ is proportional to
$S_{\rm base}/V_{\rm base}^2$, so converting an admittance from SI to p.u.\
\emph{multiplies} it by $V_{\rm base}^2/S_{\rm base}$.

%==============================================================================
\section{Conversion Formulas: SI $\to$ p.u.}
%==============================================================================

Every quantity is converted by dividing by its base:

\begin{align}
  V_{\rm pu} &= \frac{V_{\rm SI}}{V_{\rm base}}
  \label{eq:vpu}\\[4pt]
  I_{\rm pu} &= \frac{I_{\rm SI}}{I_{\rm base}}
             = I_{\rm SI} \cdot \frac{V_{\rm base}}{S_{\rm base}}
  \label{eq:ipu}\\[4pt]
  S_{\rm pu} &= \frac{S_{\rm SI}}{S_{\rm base}}
  \label{eq:spu}\\[4pt]
  Y_{\rm pu} &= \frac{Y_{\rm SI}}{Y_{\rm base}}
             = Y_{\rm SI} \cdot \frac{V_{\rm base}^2}{S_{\rm base}}
  \label{eq:ypu}\\[4pt]
  Z_{\rm pu} &= \frac{Z_{\rm SI}}{Z_{\rm base}}
             = Z_{\rm SI} \cdot \frac{S_{\rm base}}{V_{\rm base}^2}
  \label{eq:zpu}
\end{align}

Conductance $G$ and susceptance $B$ are the real and imaginary parts of $Y$,
so they follow the same scaling:

\begin{equation}
  G_{\rm pu} = G_{\rm SI} \cdot \frac{V_{\rm base}^2}{S_{\rm base}},
  \qquad
  B_{\rm pu} = B_{\rm SI} \cdot \frac{V_{\rm base}^2}{S_{\rm base}}
\end{equation}

No separate treatment is needed; they scale identically to $Y$.

%==============================================================================
\section{The Admittance Matrix $Y_{\rm bus}$: Two-Bus Scaling}
%==============================================================================

Element $(i,j)$ of $Y_{\rm bus}$ relates bus $j$ current injection to bus
$i$ voltage.  In a multi-voltage network $V_{{\rm base},i}$ and
$V_{{\rm base},j}$ can differ.  Starting from the nodal current equation
$I = Y_{\rm bus} V$ and applying per-unit conversion to both sides:

\begin{equation}
  \underbrace{\frac{I_{{\rm SI},i}}{I_{{\rm base},i}}}_{I_{{\rm pu},i}}
  =
  \sum_j
  \underbrace{\left(
    Y_{{\rm SI},ij} \cdot \frac{V_{{\rm base},j}}{I_{{\rm base},i}}
  \right)}_{Y_{{\rm pu},ij}}
  \cdot
  \underbrace{\frac{V_{{\rm SI},j}}{V_{{\rm base},j}}}_{V_{{\rm pu},j}}
\end{equation}

Substituting $I_{{\rm base},i} = S_{\rm base}/V_{{\rm base},i}$:

\begin{equation}
  \boxed{
    Y_{{\rm bus,pu},ij}
    =
    Y_{{\rm bus,SI},ij}
    \cdot
    \frac{V_{{\rm base},i} \cdot V_{{\rm base},j}}{S_{\rm base}}
  }
  \label{eq:ybuspu}
\end{equation}

The diagonal ($i=j$) reduces to:

\begin{equation}
  Y_{{\rm bus,pu},ii}
  =
  Y_{{\rm bus,SI},ii}
  \cdot
  \frac{V_{{\rm base},i}^2}{S_{\rm base}}
\end{equation}

which is exactly the scalar formula \eqref{eq:ypu} applied with the local
bus base.

In code, the full matrix conversion is a single outer-product multiply:

\begin{verbatim}
def ybus_si_to_pu(Y_si, Vbase_bus, S_base):
    scale = np.outer(Vbase_bus, Vbase_bus) / S_base
    return Y_si * scale
\end{verbatim}

\texttt{np.outer(Vbase\_bus, Vbase\_bus)} constructs the matrix
$[V_{{\rm base},i} \cdot V_{{\rm base},j}]_{ij}$ in one call.

%==============================================================================
\section{Single-Voltage vs.\ Multi-Voltage Networks}
%==============================================================================

\paragraph{Single-voltage network} (all buses share the same $V_{\rm base}$):
\begin{equation}
  Y_{{\rm pu},ij} = Y_{{\rm SI},ij} \cdot \frac{V_{\rm base}^2}{S_{\rm base}}
\end{equation}
A single scalar multiplication suffices.

\paragraph{Multi-voltage network} (e.g., Heo1: $3$--$380\,\mathrm{kV}$):
\begin{equation}
  Y_{{\rm pu},ij} = Y_{{\rm SI},ij} \cdot \frac{V_{{\rm base},i} \cdot V_{{\rm base},j}}{S_{\rm base}}
\end{equation}
An outer-product matrix must be used.  Using a single scalar $V_{\rm base}$
would incorrectly scale the off-diagonal transformer branches (e.g.,
$110\,\mathrm{kV} \leftrightarrow 20\,\mathrm{kV}$), causing the Ybus to be
wrong and all downstream power-flow residuals to be corrupted.

This is the root cause of one of the LVN bugs described in the experimental
results.

%==============================================================================
\section{Branch Admittance Entries}
%==============================================================================

The PPC dataset stores per-branch admittance quantities.  Each carries a
different scaling depending on which bus end it is associated with:

\begin{center}
\begin{tabular}{llll}
\toprule
Quantity & Associated end & Conversion factor \\
\midrule
\texttt{y\_series\_from} & from-bus self-admittance & $\times\, V_f^2 / S_{\rm base}$ \\
\texttt{y\_series\_to}   & to-bus self-admittance   & $\times\, V_t^2 / S_{\rm base}$ \\
\texttt{y\_series\_ft}   & from–to mutual term      & $\times\, V_f V_t / S_{\rm base}$ \\
\bottomrule
\end{tabular}
\end{center}

The mutual term \texttt{y\_series\_ft} uses the product of \emph{both} bus
bases, consistent with \eqref{eq:ybuspu}.  In code:

\begin{verbatim}
Branch_y_series_from *= Vf**2 / S_base
Branch_y_series_to   *= Vt**2 / S_base
Branch_y_series_ft   *= Vf * Vt / S_base
\end{verbatim}

%==============================================================================
\section{Consistency Check: Power Equation in p.u.}
%==============================================================================

A key property of the per-unit system is that the power-flow equations
retain the same algebraic form.  Starting from $S_{\rm SI} = V_{\rm SI}
\cdot \overline{I_{\rm SI}}$:

\begin{equation}
  S_{\rm pu}
  = \frac{S_{\rm SI}}{S_{\rm base}}
  = \frac{V_{\rm SI} \cdot \overline{I_{\rm SI}}}{S_{\rm base}}
  = \frac{V_{\rm SI}}{V_{\rm base}}
    \cdot \overline{\frac{I_{\rm SI}}{I_{\rm base}}}
    \cdot \underbrace{\frac{V_{\rm base} \cdot I_{\rm base}}{S_{\rm base}}}_{=\,1}
  = V_{\rm pu} \cdot \overline{I_{\rm pu}}
\end{equation}

So $S_{\rm pu} = V_{\rm pu} \cdot \overline{I_{\rm pu}}$ holds in p.u.\ with
no additional factors.  Similarly, the Newton--Raphson power-flow equation

\begin{equation}
  S_i = V_i \cdot \overline{(Y_{\rm bus}\, V)_i}
\end{equation}

holds \emph{identically} in both SI and p.u., which is the fundamental
advantage of the per-unit system: the solver code is the same regardless of
the physical scale.

%==============================================================================
\section{Changing $S_{\rm base}$ (Re-basing)}
%==============================================================================

If $V_{\rm base}$ is kept fixed and only $S_{\rm base}$ changes from
$S_1 \to S_2$:

\begin{center}
\begin{tabular}{lcc}
\toprule
Quantity & Transformation & Direction \\
\midrule
$V_{\rm pu}$   & unchanged              & — \\
$S_{\rm pu}$   & $\times\, S_1/S_2$    & same as $S$ \\
$Y_{\rm pu}$   & $\times\, S_1/S_2$    & same as $S$ \\
$I_{\rm pu}$   & $\times\, S_1/S_2$    & same as $S$ \\
$Z_{\rm pu}$   & $\times\, S_2/S_1$    & inverse of $S$ \\
\bottomrule
\end{tabular}
\end{center}

\paragraph{Practical example: LVN re-basing from 1\,MVA to 100\,MVA.}
LVN datasets (e.g., Heo1) are often generated with $S_{\rm base} =
1\,\mathrm{MVA}$.  Most HV/MV models use $100\,\mathrm{MVA}$.  Re-basing
from $1\,\mathrm{MVA}$ to $100\,\mathrm{MVA}$ ($S_2 / S_1 = 100$):

\begin{itemize}[noitemsep]
  \item $V_{\rm pu}$: \textbf{unchanged}.
  \item $S_{\rm pu}$, $Y_{\rm pu}$, $\Delta P$, $\Delta Q$:
    \textbf{divided by 100} (100$\times$ smaller).
  \item $Z_{\rm pu}$: multiplied by 100.
\end{itemize}

The physical solution is identical; only the numerical scale of the
quantities seen by the model changes.  After re-basing, the LVN residuals
$\Delta P$, $\Delta Q$ will be on the same order of magnitude as HV/MV
cases, which helps the loss landscape match what the model was trained on.

%==============================================================================
\section{Summary Table}
%==============================================================================

\begin{table}[h]
\centering
\caption{Per-unit conversion factors (SI $\to$ p.u.) for all key quantities.
  Single-voltage: one global $V_{\rm base}$.
  Multi-voltage: per-bus $V_{{\rm base},i}$.}
\label{tab:pu_summary}
\begin{tabular}{llcc}
\toprule
Quantity & Multiply by & Single-voltage & Multi-voltage \\
\midrule
$V$          & $1/V_{\rm base}$
             & per-bus scalar & per-bus scalar \\
$S$          & $1/S_{\rm base}$
             & global scalar  & global scalar \\
$I$          & $V_{\rm base}/S_{\rm base}$
             & per-bus scalar & per-bus scalar \\
$G,\,B,\,Y$ & $V_{\rm base}^2/S_{\rm base}$
             & one scalar     & per-bus scalar \\
$Y_{{\rm bus},ij}$ & $V_{{\rm base},i}\cdot V_{{\rm base},j}/S_{\rm base}$
             & one scalar     & \textbf{outer-product matrix} \\
$Z$          & $S_{\rm base}/V_{\rm base}^2$
             & one scalar     & per-bus scalar \\
\bottomrule
\end{tabular}
\end{table}

The only entry that requires a full matrix operation is
$Y_{{\rm bus},ij}$ in the multi-voltage case.  All other quantities are
element-wise operations.  The PPC code's
\texttt{np.outer(Vbase\_bus, Vbase\_bus)\,/\,S\_base} is correct for the
general case and reduces to the scalar formula when all bus bases are equal.

\end{document}
